\subsection{Bounded ordinary act cones through twelve corner elements}

\begin{lemma}[Proved: Tor persistence through algebra dimension three]
\label{CP2:A-three-dimensional-bounded-Tor}
Let $A$ be a finite-dimensional unital rational algebra with
$\dim_{\mathbb Q}A\leq3$, and let $W,V$ be finite right and left
$A$-modules. If $\operatorname{Tor}^A_j(W,V)=0$ for all sufficiently
large $j$, then $\operatorname{Tor}^A_2(W,V)=0$.
\end{lemma}
\begin{proof}
The nonlocal case is Lemma~\ref{A:nonlocal-four}; semisimple
algebras have zero positive Tor. Suppose therefore that $A$ is
local with nonzero radical $J$. Its residue division algebra
$D=A/J$ must be $\mathbb Q$. Indeed the nonzero space $J/J^2$
is a unital left $D$-module, so
$\dim A\geq\dim D+\dim(J/J^2)\geq2\dim D$.
Any residue dimension greater than one would contradict
$\dim A\leq3$.

If $J^2=0$, a minimal free cover of $V$ has semisimple kernel
$\Omega V\cong\mathbb Q^{\oplus t}$. If $V$ is not projective,
then $t>0$. If $W$ is also not projective, the minimal-resolution
argument of Lemma~\ref{A:local-target-facts} gives
\[
 \operatorname{Tor}^A_j(W,V)
 \cong \operatorname{Tor}^A_{j-1}(W,\mathbb Q)^{\oplus t}
 \ne0\qquad(j\geq2).
\]
Thus eventual vanishing forces one factor to be projective,
which kills every positive Tor group.

If $J^2\ne0$, then $\dim J=2$, $\dim J^2=1$ and $J^3=0$.
A lift $x$ of a basis of $J/J^2$ generates $J$ as an ideal,
and $x^2$ spans $J^2$. Therefore $A\cong\mathbb Q[x]/(x^3)$.
Jordan decomposition expresses either module as a direct sum
of $A/(x^a)$, $1\leq a\leq3$. For two nonfree summands
with $a,b<3$, the two-periodic free resolution, with maps
$x^a$ and $x^{3-a}$, gives in every positive degree the
nonzero Tor dimension
\[
 \min(a,b)+\min(3-a,b)-b
 =\min\{a,b,3-a,3-b\}>0.
\]
Free summands give zero positive Tor. Eventual vanishing again
implies the required degree-two vanishing.
\end{proof}

\begin{lemma}[Proved: the full-column projective orbit kernel]
\label{CP2:A-full-column-act-kernel}
Let $M$ have rectangular minimal band $K=I\times J$, and let
$S$ be a finite nonempty right act. Suppose that every core
orbit in $SK$ has exactly $|J|$ states. If there are $o$ such
orbits, the constant-action quotient $S\to\bar S$ of
Lemma~\ref{CP2:A-constant-quotient-H2} gives an exact sequence
of right $\mathbb ZM$-modules
\[
 0\longrightarrow(e\mathbb ZM)^{\oplus(o-1)}
 \longrightarrow\widetilde{\mathbb ZS}
 \longrightarrow\widetilde{\mathbb Z\bar S}
 \longrightarrow0,
\]
where $e$ is a minimal idempotent. The induced act-cone map is
an integral homology isomorphism in every degree at least three.
\end{lemma}
\begin{proof}
Label each core orbit by $s_{a,j}$, $j\in J$. Full size makes
this labelling bijective. Every orbit is individually invariant
under $M$, and
\[
 s_{a,j}m=s_{a,\tau_m(j)}.
\]
The generating constant-action relations identify precisely
the states with the same column label in different core orbits.
They identify no transient states, and these relations are
already invariant under every outside actor. The kernel thus
has the integral basis
\[
 [s_{a,j}]-[s_{0,j}],\qquad 1\leq a<o,\quad j\in J.
\]
For each $a$ this is the right ideal module $e\mathbb ZM$,
which is projective. The reduced act-cone complex is
$(\widetilde{\mathbb ZS}\otimes_{\mathbb ZM}^{\mathbb L}
\mathbb Z)[1]$. Its long exact Tor sequence gives the asserted
comparison in degrees at least three. In particular it passes
bounded rational homology and preserves rational fourth homology.
\end{proof}

\begin{proposition}[Proved: three columns and at most three outside elements]
\label{CP2:A-three-by-three-three-outside-bounded-acts}
Let $M$ have minimal band $3\times3$ and at most three outside
elements. Every finite nonempty ordinary right-act cone with
bounded rational homology has zero rational degree-two products.
There is no restriction on its number of states or second Betti number.
\end{proposition}
\begin{proof}
Take the actual constant-action quotient $S\to\bar S$.
Its pullback on rational $H^2$ is surjective by
Lemma~\ref{CP2:A-constant-quotient-H2}. Write $c_j\in\bar S$
for the constant value of a band element in old column $j$,
and put $T=\{c_j:j\in\{0,1,2\}\}$.

If $|T|=1$, its state $c$ is fixed by the whole minimal ideal
and hence by $M$: for a band idempotent $e$,
$cm=c(em)=c$. This fixed state gives a section of the action
bar projection to $BM$. Every second cohomology class of the
quotient cone is relative to the monoid vertex, so its products
vanish. Surjectivity of the actual $H^2$ pullback finishes.

If $|T|=2$, collapse old columns with equal values $c_j$,
keeping every outside element distinct. The identity
$c_jm=c_{\tau_m(j)}$ shows that the column equivalence is
preserved by every right actor; left multiplication leaves
columns unchanged. It is therefore an actual monoid congruence.
The quotient $M'$ has minimal band $3\times2$ and the same
at most three outside elements, and the action on $\bar S$
factors through $M'$. The actual act-cone functor gives an
$H^2$ isomorphism by Lemma~\ref{CP2:A-column-quotient-act-H2}.
Every ordinary act cone over $M'$ has zero rational second
products by Proposition~\ref{CP2:A-three-row-five-outside-acts}.
Pullback through both functors proves the claim in this case.
No boundedness of the quotient was assumed or inferred.

Finally suppose $|T|=3$. An equality between two column states
in any original core orbit would identify the corresponding
$c_j$ in $\bar S$. Thus all original core orbits are full
three-state orbits. The preceding projective-kernel lemma
makes the quotient cone bounded in rational homology and
identifies its fourth homology with that of the original cone.
Constant-action compression
(Theorem~\ref{CP2:A-constant-act-compression}) identifies the
quotient cone, by an actual homotopy equivalence, with a monoid
$P$ having band $\{0,1,2\}\times\bar S$ and the same contracted
algebra
\[
 A=\mathbb QM/\mathbb QK,\qquad \dim A\leq3.
\]
The band--Tor formula expresses $H_{j+2}(BP;\mathbb Q)$ as
$\operatorname{Tor}^A_j(W,V)$ for its column and row augmentation
modules. Boundedness and Lemma~\ref{CP2:A-three-dimensional-bounded-Tor}
give $H_4(BP;\mathbb Q)=0$. The original fourth homology is
therefore zero as well. This proves all second-product vanishings.
\end{proof}

\begin{theorem}[Proved: bounded act cones through order twelve]
\label{CP2:A-bounded-twelve-acts}
For every finite monoid $M$ of order at most twelve and every
finite nonempty right act $S$, bounded rational homology of
$B\mathcal C(M,S)$ forces all rational degree-two products
to vanish. The left-act version holds by opposition.
\end{theorem}
\begin{proof}
Orders through eleven are Theorem~\ref{CP2:A-bounded-eleven-acts}.
The same initial reductions apply at order twelve: if
$H^2(BM;\mathbb Q)=0$, every second class is relative to the
monoid vertex; at most two minimal rows are covered by
Lemma~\ref{CP2:A-high-two-row-act}. Otherwise a possible
nonzero product requires at least three rows and two columns.
A nontrivial minimal group would then require at least
$3\cdot2\cdot2+1=13$ elements. Thus the minimal ideal is a band.
Writing $q=|M\setminus K|$, the complete list is
\[
 (r,c,q)=(3,2,q\leq6),\quad(4,2,q\leq4),\quad
 (5,2,q\leq2),\quad(3,3,q\leq3).
\]
The first family is Proposition~\ref{CP2:A-three-row-six-outside-acts},
whose complete27-table projective-row calculation has already
been independently checked. The next two use
Lemma~\ref{CP2:A-two-column-four-outside-bounded}.
The last family is the preceding proposition.
\end{proof}

\begin{corollary}[Proved: two nineteen-arrow ordinary/projective budgets]
\label{CP2:A-nineteen-twelve-small-opposite}
No two-object category in either oriented budget
$(12,1,6)$ or $(12,2,5)$ has classifying space of type $\CP^2$.
More generally, if its source corner has order at most twelve
and its opposite corner has projective trivial right rational
module, bounded rational homology forces all degree-two products
to vanish, for every nonempty connecting biset.
\end{corollary}
\begin{proof}
Exact coinvariants at the opposite corner are the actual
finite orbit quotient of the connecting biset. Filtering the
two-sided bar by that corner identifies the actual quotient
functor to the ordinary source-act cone as a rational homology
and cohomology-ring equivalence. In particular it preserves
boundedness. Apply the theorem. Both monoids of order two
have semisimple rational monoid algebras (the group $C_2$ and
the identity-adjoined singleton idempotent), so the required
projectivity holds for both displayed budgets.
\end{proof}

\begin{remark*}[Precise new scope]
This extension reconstructs the full-column projective kernel
and the dimension-three Tor argument. It does not use the
inherited, stronger contracted-four assertion or its absent
certificates. The only reused finite source calculation is the
already adopted27-table three-row/six-outside gate. No general
boundedness transfer through a faithful actor-image quotient is
asserted. Existing nineteen-arrow inventories and completed
budget snapshots are unchanged until a separate exact scope join.
\end{remark*}
